\section{Multiplicative valuations}

\begin{definition}[Multiplicative valuation associated to a discrete valuation]
    Let \(K\) be a field, and let
    \[
        v:K^\times\longrightarrow \mathbb{Z}
    \]
    be a normalized discrete valuation. We extend \(v\) to \(K\) by setting
    \[
        v(0)=+\infty.
    \]
    Fix a real number
    \[
        0<\gamma<1.
    \]
    Define a function
    \[
        \varphi:K\longrightarrow \mathbb{R}_{\geq 0}
    \]
    by
    \[
        \varphi(a)=
        \begin{cases}
            \gamma^{v(a)}, & a\neq 0,\\
            0, & a=0.
        \end{cases}
    \]
    Then \(\varphi\) is called the multiplicative valuation associated to \(v\)
    with parameter \(\gamma\).
\end{definition}

\begin{proposition}
    The function
    \[
        \varphi:K\longrightarrow \mathbb{R}_{\geq 0}
    \]
    defined above satisfies the following properties:
    \begin{enumerate}
        \item \(\varphi(a)=0\) if and only if \(a=0\).
        \item For all \(a,b\in K\),
              \[
                  \varphi(ab)=\varphi(a)\varphi(b).
              \]
        \item For all \(a,b\in K\),
              \[
                  \varphi(a+b)\leq \max\{\varphi(a),\varphi(b)\}.
              \]
    \end{enumerate}
\end{proposition}

\begin{proof}
    The first assertion follows immediately from the definition, because
    \[
        \gamma^{v(a)}>0
    \]
    whenever \(a\neq 0\), while \(\varphi(0)=0\).

    For the multiplicativity, if \(a,b\neq 0\), then
    \[
        v(ab)=v(a)+v(b).
    \]
    Hence
    \[
        \varphi(ab)
        =
        \gamma^{v(ab)}
        =
        \gamma^{v(a)+v(b)}
        =
        \gamma^{v(a)}\gamma^{v(b)}
        =
        \varphi(a)\varphi(b).
    \]
    If \(a=0\) or \(b=0\), both sides are zero, so the formula still holds.

    Finally, assume first that \(a+b\neq 0\). Since \(v\) is a valuation,
    \[
        v(a+b)\geq \min\{v(a),v(b)\}.
    \]
    Because \(0<\gamma<1\), the function \(x\mapsto \gamma^x\) is decreasing.
    Therefore
    \[
        \gamma^{v(a+b)}
        \leq
        \gamma^{\min\{v(a),v(b)\}}
        =
        \max\{\gamma^{v(a)},\gamma^{v(b)}\}.
    \]
    Thus
    \[
        \varphi(a+b)\leq \max\{\varphi(a),\varphi(b)\}.
    \]
    If \(a+b=0\), then \(\varphi(a+b)=0\), and the same inequality is obvious.
\end{proof}

\begin{definition}[Multiplicative valuation]
    Let \(K\) be a field. A function
    \[
        \varphi:K\longrightarrow \mathbb{R}_{\geq 0}
    \]
    is called a multiplicative valuation, or a non-archimedean absolute value, if
    it satisfies:
    \begin{enumerate}
        \item \(\varphi(a)=0\) if and only if \(a=0\);
        \item
              \[
                  \varphi(ab)=\varphi(a)\varphi(b)
              \]
              for all \(a,b\in K\);
        \item
              \[
                  \varphi(a+b)\leq \max\{\varphi(a),\varphi(b)\}
              \]
              for all \(a,b\in K\).
    \end{enumerate}
\end{definition}

\begin{proposition}[Basic properties of multiplicative valuations]
    Let
    \[
        \varphi:K\longrightarrow \mathbb{R}_{\geq 0}
    \]
    be a multiplicative valuation. Then for all \(a,b\in K^\times\), the following
    hold:
    \begin{enumerate}
        \item
              \[
                  \varphi(1)=1,
                  \qquad
                  \varphi(-1)=1.
              \]
        \item
              \[
                  \varphi(-a)=\varphi(a).
              \]
        \item
              \[
                  \varphi(a^{-1})=\varphi(a)^{-1}.
              \]
        \item If
              \[
                  \varphi(a)<\varphi(b),
              \]
              then
              \[
                  \varphi(a+b)=\varphi(b).
              \]
    \end{enumerate}
\end{proposition}

\begin{proof}
    Since \(1\cdot 1=1\), we have
    \[
        \varphi(1)=\varphi(1)^2.
    \]
    Since \(\varphi(1)\neq 0\), it follows that
    \[
        \varphi(1)=1.
    \]
    Also, since \((-1)^2=1\), we get
    \[
        \varphi(-1)^2=\varphi(1)=1.
    \]
    Since \(\varphi(-1)>0\), we have
    \[
        \varphi(-1)=1.
    \]

    Hence, for \(a\in K^\times\),
    \[
        \varphi(-a)=\varphi((-1)a)=\varphi(-1)\varphi(a)=\varphi(a).
    \]

    Also,
    \[
        1=\varphi(1)=\varphi(aa^{-1})=\varphi(a)\varphi(a^{-1}),
    \]
    so
    \[
        \varphi(a^{-1})=\varphi(a)^{-1}.
    \]

    Finally, suppose that
    \[
        \varphi(a)<\varphi(b).
    \]
    By the non-archimedean triangle inequality,
    \[
        \varphi(a+b)\leq \max\{\varphi(a),\varphi(b)\}=\varphi(b).
    \]
    On the other hand,
    \[
        b=(a+b)-a.
    \]
    Hence
    \[
        \varphi(b)
        \leq
        \max\{\varphi(a+b),\varphi(a)\}.
    \]
    If \(\varphi(a+b)<\varphi(b)\), then both \(\varphi(a+b)\) and \(\varphi(a)\) are
    strictly smaller than \(\varphi(b)\), contradicting the above inequality. Therefore
    \[
        \varphi(a+b)=\varphi(b).
    \]
\end{proof}

\begin{definition}[Valuation ring and valuation ideal]
    Let \(\varphi\) be a multiplicative valuation on \(K\). Define
    \[
        R_\varphi
        :=
        \{a\in K\mid \varphi(a)\leq 1\},
    \]
    and
    \[
        P_\varphi
        :=
        \{a\in K\mid \varphi(a)<1\}.
    \]
    Then \(R_\varphi\) is called the valuation ring of \(\varphi\), and
    \(P_\varphi\) is called the valuation ideal of \(\varphi\).
\end{definition}

\begin{proposition}
    Let \(\varphi\) be a multiplicative valuation on \(K\). Then \(R_\varphi\) is a
    subring of \(K\), and \(P_\varphi\) is an ideal of \(R_\varphi\).
\end{proposition}

\begin{proof}
    We first show that \(R_\varphi\) is a subring. Clearly,
    \[
        \varphi(0)=0\leq 1,
        \qquad
        \varphi(1)=1,
    \]
    so \(0,1\in R_\varphi\). If \(a,b\in R_\varphi\), then
    \[
        \varphi(ab)=\varphi(a)\varphi(b)\leq 1,
    \]
    hence \(ab\in R_\varphi\). Also,
    \[
        \varphi(-a)=\varphi(a)\leq 1,
    \]
    so \(-a\in R_\varphi\). Finally,
    \[
        \varphi(a+b)\leq \max\{\varphi(a),\varphi(b)\}\leq 1,
    \]
    so \(a+b\in R_\varphi\). Therefore \(R_\varphi\) is a subring of \(K\).

    Next we show that \(P_\varphi\) is an ideal of \(R_\varphi\). If \(a,b\in P_\varphi\),
    then
    \[
        \varphi(a+b)\leq \max\{\varphi(a),\varphi(b)\}<1,
    \]
    so \(a+b\in P_\varphi\). Also, if \(r\in R_\varphi\) and \(a\in P_\varphi\), then
    \[
        \varphi(ra)=\varphi(r)\varphi(a)<1.
    \]
    Hence \(ra\in P_\varphi\). Thus \(P_\varphi\) is an ideal of \(R_\varphi\).
\end{proof}

\begin{remark}
    If \(\varphi(a)=\gamma^{v(a)}\) for a normalized discrete valuation \(v\), where
    \(0<\gamma<1\), then
    \[
        R_\varphi
        =
        \{a\in K\mid \varphi(a)\leq 1\}
        =
        \{a\in K\mid v(a)\geq 0\}
        =
        R_v,
    \]
    and
    \[
        P_\varphi
        =
        \{a\in K\mid \varphi(a)<1\}
        =
        \{a\in K\mid v(a)>0\}
        =
        P_v.
    \]
\end{remark}

\begin{proposition}[From multiplicative valuations to additive valuations]
    Let \(\varphi\) be a multiplicative valuation on \(K\), and fix
    \[
        0<\gamma<1.
    \]
    Define
    \[
        v_\varphi:K^\times\longrightarrow \mathbb{R}
    \]
    by
    \[
        v_\varphi(a):=\log_\gamma \varphi(a).
    \]
    Then \(v_\varphi\) is an additive valuation on \(K\), that is,
    \[
        v_\varphi(ab)=v_\varphi(a)+v_\varphi(b),
    \]
    and
    \[
        v_\varphi(a+b)\geq \min\{v_\varphi(a),v_\varphi(b)\}
    \]
    whenever \(a+b\neq 0\).
\end{proposition}

\begin{proof}
    For \(a,b\in K^\times\), we have
    \[
        \varphi(ab)=\varphi(a)\varphi(b).
    \]
    Taking logarithms with base \(\gamma\), we get
    \[
        v_\varphi(ab)
        =
        \log_\gamma \varphi(ab)
        =
        \log_\gamma \varphi(a)+\log_\gamma \varphi(b)
        =
        v_\varphi(a)+v_\varphi(b).
    \]

    Now suppose \(a+b\neq 0\). Since \(\varphi\) is non-archimedean,
    \[
        \varphi(a+b)\leq \max\{\varphi(a),\varphi(b)\}.
    \]
    Because \(0<\gamma<1\), the function \(\log_\gamma(x)\) is decreasing. Therefore
    \[
        \log_\gamma \varphi(a+b)
        \geq
        \log_\gamma \max\{\varphi(a),\varphi(b)\}.
    \]
    But
    \[
        \log_\gamma \max\{\varphi(a),\varphi(b)\}
        =
        \min\{\log_\gamma \varphi(a),\log_\gamma \varphi(b)\}.
    \]
    Hence
    \[
        v_\varphi(a+b)\geq \min\{v_\varphi(a),v_\varphi(b)\}.
    \]
    Thus \(v_\varphi\) is an additive valuation.
\end{proof}

\begin{definition}[Value group]
    Let \(\varphi\) be a multiplicative valuation on \(K\). Then
    \[
        \varphi(K^\times)\subseteq \mathbb{R}_{>0}
    \]
    is a multiplicative subgroup of \(\mathbb{R}_{>0}\). It is called the value group
    of \(\varphi\).
\end{definition}

\begin{definition}[Equivalent multiplicative valuations]
    Let \(K\) be a field, and let
    \[
        \varphi,\psi:K\longrightarrow \mathbb{R}_{\geq 0}
    \]
    be two multiplicative valuations on \(K\). We say that \(\varphi\) and \(\psi\)
    are equivalent if there exists a positive real number
    \[
        \gamma\in \mathbb{R}_{>0}
    \]
    such that
    \[
        \psi(a)=\varphi(a)^{\gamma}
        \qquad
        \text{for all } a\in K.
    \]
\end{definition}

\begin{definition}[Metric induced by a multiplicative valuation]
    Let \(\varphi\) be a multiplicative valuation on a field \(K\). Define
    \[
        d:K\times K\longrightarrow \mathbb{R}_{\geq 0}
    \]
    by
    \[
        d(a,b)=\varphi(a-b).
    \]
    Then \(d\) is a metric on \(K\).
\end{definition}

\begin{proposition}
    Let \(K\) be a field equipped with a multiplicative valuation \(\varphi\).
    Then the function
    \[
        d(a,b)=\varphi(a-b)
    \]
    defines a metric on \(K\). In particular, \(K\) is a Hausdorff topological space
    with respect to the topology induced by \(d\).
\end{proposition}

\begin{proof}
    We verify the axioms of a metric.

    First, for any \(a,b\in K\), we have
    \[
        d(a,b)=\varphi(a-b)\geq 0.
    \]
    Moreover,
    \[
        d(a,b)=0
        \Longleftrightarrow
        \varphi(a-b)=0
        \Longleftrightarrow
        a-b=0
        \Longleftrightarrow
        a=b.
    \]

    Next, since \(\varphi(-1)=1\), we have
    \[
        d(a,b)=\varphi(a-b)
        =
        \varphi\bigl(-(b-a)\bigr)
        =
        \varphi(b-a)
        =
        d(b,a).
    \]

    Finally, for any \(a,b,c\in K\), the valuation inequality gives
    \[
        d(a,c)
        =
        \varphi(a-c)
        =
        \varphi\bigl((a-b)+(b-c)\bigr)
        \leq
        \varphi(a-b)+\varphi(b-c).
    \]
    Hence
    \[
        d(a,c)\leq d(a,b)+d(b,c).
    \]

    Therefore \(d\) is a metric on \(K\). Since every metric space is Hausdorff,
    the field \(K\) is Hausdorff with respect to the topology induced by \(d\).
\end{proof}
